% ======================================================================
% 第 11 章：重整化与对称性
% （原书第 347-391 页；PDF 第 370-415 页）
% ======================================================================
\chapter{重整化与对称性}

既然我们已经确定了量子场论中紫外发散的普遍结构，似乎很自然地应该继续研究这些发散在费曼图计算中的含义。然而，我们现在将把这一问题搁置到第~12 章，转而讨论一个看似毫不相关的方向。在第~8 章和第~9.3 节中，我们注意到了量子场论与统计力学之间的形式联系。标量场论最接近的形式类比被认为是铁磁体或其他允许二级相变的系统的连续描述。这一类比提出了这样一种可能性：在量子场论中，场同样可能取一个非零的整体值。与磁体一样，这一整体场可能具有方向性特征，从而破坏拉格朗日量的某种对称性。在这种情况下，我们说该场论具有隐藏对称性或\emph{自发对称性破缺}。我们将在本章致力于分析这一对称性破坏机制。

自发对称性破缺是量子场论研究中的一个核心概念，原因有二。首先，它在量子场论对自然界的应用中扮演着重要角色。在本书中，我们将看到此类应用的两个截然不同的例子：第~13 章将把隐藏对称性理论应用于统计力学，具体而言是二级相变附近热力学变量的行为。稍后，在第~20 章，我们将看到隐藏对称性是弱相互作用理论中的一个基本要素。自发对称性破缺在强相互作用理论以及寻找基本物理统一模型中也都有应用。

但从理论角度来看，自发对称性破缺也同样引人入胜。具有自发对称性破缺的量子场论包含紫外发散。因此，很自然地要问：这些发散是否受到理论基本对称性的约束？对这一问题的回答，最早由 Benjamin Lee 给出，%
\footnote{Lee 的分析有一个精彩的总结，见他的讲义：B.~Lee, \emph{Chiral Dynamics} (Gordon and Breach, New York, 1972)。}
将使我们进一步洞察紫外发散的本质以及重整化的意义。

\section{自发对称性破缺}

我们从经典场论中自发对称性破缺的分析开始。首先考虑熟悉的 $\phi^4$ 理论拉格朗日量，
\begin{equation}
  \mathcal{L} = \tfrac{1}{2}(\partial_\mu\phi)^2 - \tfrac{1}{2}m^2\phi^2
    - \frac{\lambda}{4!}\phi^4 ,
\end{equation}
但把 $m^2$ 替换为一个负参数 $-\mu^2$：
\begin{equation}
  \mathcal{L} = \tfrac{1}{2}(\partial_\mu\phi)^2 + \tfrac{1}{2}\mu^2\phi^2
    - \frac{\lambda}{4!}\phi^4 .
  \label{eq:11.1}
\end{equation}
该拉格朗日量具有一个离散对称性：它在运算 $\phi\to-\phi$ 下不变。相应的哈密顿量为
\begin{equation}
  H = \int d^3x\ \Bigl[ \tfrac{1}{2}\pi^2 + \tfrac{1}{2}(\nabla\phi)^2
    + V(\phi) \Bigr] .
\end{equation}
最小能量经典组态是均匀场 $\phi(x) = \phi_0$，其中 $\phi_0$ 的选择要使势
\begin{equation}
  V(\phi) = -\frac{1}{2}\mu^2\phi^2 + \frac{\lambda}{4!}\phi^4
\end{equation}
取极小值（见图~\ref{fig:11.1}）。该势有两个极小值，由下式给出
\begin{equation}
  \phi_0 = \pm v = \pm \sqrt{\frac{6\mu^2}{\lambda}} .
  \label{eq:11.2}
\end{equation}
常数 $v$ 称为 $\phi$ 的\emph{真空期望值}。

\begin{figure}[t]
  \centering
  \includegraphics[width=0.45\textwidth]{fig11.1.png}
  \caption{离散情形下自发对称性破缺的势。}
  \label{fig:11.1}
\end{figure}

为解释这一理论，假设系统处于其中一个极小值（比如说正的那个）附近。此时定义
\begin{equation}
  \phi(x) = v + \sigma(x),
  \label{eq:11.3}
\end{equation}
并用 $\sigma(x)$ 重新写出 $\mathcal{L}$ 是方便的。把 \eqref{eq:11.3} 代入 \eqref{eq:11.1}，我们发现线性于 $\sigma$ 的项为零（这是必然的，因为势的极小值位于 $\sigma = 0$ 处）。同时略去常数项，我们得到拉格朗日量
\begin{equation}
  \mathcal{L} = \tfrac{1}{2}(\partial_\mu\sigma)^2
    - \tfrac{1}{2}(2\mu^2)\sigma^2 - \sqrt{\frac{\lambda}{6}}\,\mu\,\sigma^3
    - \frac{\lambda}{4!}\sigma^4 .
  \label{eq:11.4}
\end{equation}
该拉格朗日量描述一个质量为 $\sqrt{2}\mu$ 的简单标量场，带有 $\sigma^3$ 和 $\sigma^4$ 相互作用。对称性 $\phi\to-\phi$ 不再明显；它的唯一体现是 \eqref{eq:11.4} 中三个系数之间的关系，这些系数以一种特殊方式仅依赖于两个参数。这是自发对称性破缺的最简单例子。

\subsection{线性西格马模型}

当被破缺的对称性是连续的而非离散时，会出现更有趣的理论。最重要的例子是前述理论的一种推广，称为\emph{线性西格马模型}，我们在习题~4.3 中曾简要考虑过它。我们将在本章中详细研究这一模型。

线性西格马模型的拉格朗日量涉及一组 $N$ 个实标量场 $\phi^i(x)$，
\begin{equation}
  \mathcal{L} = \tfrac{1}{2}(\partial_\mu\phi^i)^2 + \tfrac{1}{2}\mu^2(\phi^i)^2
    - \frac{\lambda}{4}\bigl[ (\phi^i)^2 \bigr]^2 ,
  \label{eq:11.5}
\end{equation}
其中在每个因子 $((\phi^i)^2)$ 中对 $i$ 求和（省略了求和号）。注意，我们把 $\phi^4$ 理论拉格朗日量中的耦合常数 $\lambda$ 重新标度了，以消除上面分析中别扭的因子 6。拉格朗日量 \eqref{eq:11.5} 在对称性
\begin{equation}
  \phi^i \to R^{ij}\phi^j ,
  \label{eq:11.6}
\end{equation}
下不变，其中 $R$ 是任意 $N\times N$ 正交矩阵。变换 \eqref{eq:11.6} 的群正是 $N$ 维旋转群，也称为 $N$ 维正交群，或简称为 $O(N)$。

同样，最低能量经典组态是常数场 $\phi^i_0$，其取值选择为使势
\begin{equation}
  V(\phi) = -\frac{1}{2}\mu^2(\phi^i)^2 + \frac{\lambda}{4}\bigl[
    (\phi^i)^2 \bigr]^2
\end{equation}
取极小值（见图~\ref{fig:11.2}）。对任何满足
\begin{equation}
  (\phi^i_0)^2 = \frac{\mu^2}{\lambda} .
\end{equation}
的 $\phi^i_0$，该势都被极小化。这一条件只确定了矢量 $\phi_0$ 的长度；其方向是任意的。习惯上选择坐标使得 $\phi^i_0$ 指向第 $N$ 个方向：
\begin{equation}
  \phi^i_0 = (0, 0, \ldots, 0, v),
  \qquad\text{其中 } v = \sqrt{\frac{\mu^2}{\lambda}} .
  \label{eq:11.7}
\end{equation}
我们现在可以通过写下
\begin{equation}
  \phi^i(x) = \bigl( \pi^k(x), v + \sigma(x) \bigr),
  \qquad k = 1, \ldots, N-1 .
  \label{eq:11.8}
\end{equation}
来定义一组平移后的场。（与习题~4.3 一样，这种记号来源于在 $N = 4$ 情形下把这一形式体系应用于 $\pi$ 介子。）

\begin{figure}[t]
  \centering
  \includegraphics[width=0.45\textwidth]{fig11.2.png}
  \caption{连续 $O(N)$ 对称性自发破缺的势，图为 $N = 2$ 的情形。沿势谷的振荡对应于无质量 $\pi$ 场。}
  \label{fig:11.2}
\end{figure}

现在直接用 $\pi$ 场和 $\sigma$ 场重写拉格朗日量 \eqref{eq:11.5} 是直接的。结果是
\begin{equation}
  \mathcal{L} = \tfrac{1}{2}(\partial_\mu\pi^k)^2 + \tfrac{1}{2}
    (\partial_\mu\sigma)^2 - \frac{1}{2}(2\mu^2)\sigma^2
    - \sqrt{\lambda}\,\mu\,\sigma^3 - \sqrt{\lambda}\,\mu\,(\pi^k)^2\sigma
    - \frac{\lambda}{4}\sigma^4 - \frac{\lambda}{2}(\pi^k)^2\sigma^2
    - \frac{\lambda}{4}\bigl[ (\pi^k)^2 \bigr]^2 .
  \label{eq:11.9}
\end{equation}
我们像在 \eqref{eq:11.4} 中那样得到一个有质量的 $\sigma$ 场，以及一组 $N-1$ 个无质量 $\pi$ 场。原来的 $O(N)$ 对称性是隐藏的，只剩下子群 $O(N-1)$，它在 $\pi$ 场之间进行旋转。参见图~\ref{fig:11.2}，我们注意到有质量 $\sigma$ 场描述 $\phi^i$ 在径向的振荡，在该方向上势的二阶导数不为零。无质量 $\pi$ 场描述 $\phi^i$ 在切向的振荡，即沿势的谷底。谷底是一个 $(N-1)$ 维曲面，所有 $N-1$ 个方向都是等价的，这反映了未被破缺的 $O(N-1)$ 对称性。

\subsection{戈德斯通定理}

当连续对称性被自发破缺时无质量粒子的出现是一个普遍结果，称为\emph{戈德斯通定理}。为精确表述这一定理，我们必须计数线性无关的连续对称变换的数目。在线性西格马模型中，对于 $N = 1$ 没有连续对称性，而对于 $N = 2$ 只有一个旋转方向。$N$ 维中的旋转可以发生在 $N(N-1)/2$ 个平面中的任意一个，因此 $O(N)$ 对称理论有 $N(N-1)/2$ 个连续对称性。在自发对称性破缺之后，剩下 $(N-1)(N-2)/2$ 个对称性，对应于 $(N-1)$ 个 $\pi$ 场的旋转。被破缺的对称性数目是二者的差，即 $N-1$。

戈德斯通定理指出，对于每一个被自发破缺的连续对称性，理论必须包含一个无质量粒子。%
\footnote{J.~Goldstone, \emph{Nuovo Cim.} \textbf{19}, 154 (1961)。J.~Goldstone, A.~Salam 和 S.~Weinberg 的一篇颇具启发性的四页论文，\emph{Phys.\ Rev.} \textbf{127}, 965 (1962)，给出了该定理的三个不同证明。}
我们刚刚看到，这一定理在线性西格马模型中成立，至少在经典层面如此。通过自发对称性破缺产生的无质量场称为\emph{戈德斯通玻色子}。物理学中看到的许多轻玻色子，例如 $\pi$ 介子，都可以（至少近似地）解释为戈德斯通玻色子。我们以经典标量场理论中戈德斯通定理的一般证明来结束本节。本章的其余部分致力于对具有隐藏对称性的理论进行量子力学分析。到本章结束时我们将看到，戈德斯通玻色子不可能从任何阶的量子修正中获得质量。

考虑一个涉及多个场 $\phi^a(x)$ 的理论，其拉格朗日量具有如下形式
\begin{equation}
  \mathcal{L} = (\text{含导数的项}) - V(\phi) .
  \label{eq:11.10}
\end{equation}
设 $\phi_0$ 是使 $V$ 取极小值的常数场，从而
\begin{equation}
  \frac{\partial V}{\partial\phi^a}\Big|_{\phi_0} = 0 .
\end{equation}
把 $V$ 在这一点附近展开，我们发现
\begin{equation}
  V(\phi) = V(\phi_0) + \frac{1}{2}(\phi-\phi_0)^a\,
    \frac{\partial^2 V}{\partial\phi^a\partial\phi^b}\Big|_{\phi_0}\,
    (\phi-\phi_0)^b + \cdots .
\end{equation}
二次项的系数
\begin{equation}
  M^2_{ab} = \frac{\partial^2 V}{\partial\phi^a\partial\phi^b}\Big|_{\phi_0},
  \label{eq:11.11}
\end{equation}
是一个对称矩阵，其特征值给出各场的质量。这些特征值不能为负，因为 $\phi_0$ 是一个极小值。为证明戈德斯通定理，我们必须证明：拉格朗日量 \eqref{eq:11.10} 的每一个不是 $\phi_0$ 的对称性的连续对称性，都会导致这个质量矩阵出现一个零特征值。

一般连续对称变换具有如下形式
\begin{equation}
  \phi^a \to \phi^a + \alpha\Delta^a(\phi),
  \label{eq:11.12}
\end{equation}
其中 $\alpha$ 是无穷小参数，$\Delta^a$ 是所有 $\phi$ 的某个函数。特殊化到常数场；此时 $\mathcal{L}$ 中的导数项消失，仅势本身必须在 \eqref{eq:11.12} 下不变。这一条件可以写为
\begin{equation}
  V(\phi^a) = V\bigl( \phi^a + \alpha\Delta^a(\phi) \bigr)
  \qquad\text{或}\qquad
  \Delta^a(\phi)\frac{\partial V}{\partial\phi^a} = 0 .
\end{equation}
现在对 $\phi^b$ 求导，并令 $\phi = \phi_0$：
\begin{equation}
  0 = \Bigl( \frac{\partial\Delta^a}{\partial\phi^b} \frac{\partial V}
    {\partial\phi^a} + \Delta^a \frac{\partial^2 V}{\partial\phi^a
      \partial\phi^b} \Bigr)\Big|_{\phi_0}
  = \Delta^a(\phi_0)\,M^2_{ab} .
  \label{eq:11.13}
\end{equation}
第一项因 $\phi_0$ 是 $V$ 的极小值而为零，所以第二项也必然为零。如果变换使 $\phi_0$ 保持不变（即，如果对称性被基态所尊重），那么 $\Delta^a(\phi_0) = 0$，这一关系是平凡的。自发破缺的对称性恰恰是 $\Delta^a(\phi_0) \neq 0$ 的那种；在这种情况下，$\Delta^a(\phi_0)$ 正是我们所要的特征值为零的矢量，因此戈德斯通定理得证。

\section{重整化与对称性：一个显式例子}

现在让我们研究具有自发对称性破缺的理论的量子力学。我们仍将以线性西格马模型为例。该理论的拉格朗日量，用平移后的场写出，由式~\eqref{eq:11.9} 给出。从这个表达式，我们可以读出费曼规则；这些规则如图~\ref{fig:11.3} 所示。

\begin{figure}[t]
  \centering
  \includegraphics[width=0.6\textwidth]{fig11.3.png}
  \caption{以平移后的场表示的线性西格马模型的费曼规则。}
  \label{fig:11.3}
\end{figure}

利用这些费曼规则，我们可以毫无困难地计算树图级振幅。然而，带圈的图通常会发散。对于有 $N_e$ 条外腿的振幅，表面发散度为
\begin{equation}
  D = 4 - N_e ,
\end{equation}
这与第~10.2 节对 $\phi^4$ 理论的讨论完全类似。（含三点顶点的图会比这个表达式所表示的更不发散，因为这个顶点有一个具有质量量纲的系数。）然而，对称性对振幅的约束比前面的讨论要弱得多：该理论有三个耦合常数（$\mu$、$\lambda$ 和 $v$），它们之间的关系只在树图级被固定，理论的发散可能与 $O(N)$ 对称性一致，也可能不一致。

发散振幅如图~\ref{fig:11.4} 所示。真空能、单点 $\sigma$ 振幅（蝌蚪图）、两点函数和四点函数中都存在发散。如果理论要可重整化，所有这些发散都必须被如下形式的抵消项所抵消
\begin{equation}
  \delta\mathcal{L} = \tfrac{1}{2}\delta_Z(\partial_\mu\phi^i)^2
    + \tfrac{1}{2}\delta_\mu\phi^i\phi^i
    + \delta_\lambda\bigl( \phi^i\phi^i \bigr)^2
    + \delta_v\phi^N ,
  \label{eq:11.15}
\end{equation}
其中只包含三个可调参数 $\delta_Z$、$\delta_\mu$ 和 $\delta_\lambda$，外加蝌蚪图抵消项 $\delta_v$。这三个抵消项是否足以抵消图~\ref{fig:11.4} 中所有的无穷大，这一点完全不显然。这正是本章开头提出的问题：具有自发对称性破缺的理论的发散是否受到基本 $O(N)$ 对称性的约束？

\begin{figure}[t]
  \centering
  \includegraphics[width=0.6\textwidth]{fig11.4.png}
  \caption{线性西格马模型中的发散振幅。}
  \label{fig:11.4}
\end{figure}

\begin{figure}[t]
  \centering
  \includegraphics[width=0.65\textwidth]{fig11.5.png}
  \caption{线性西格马模型中抵消项顶点的费曼规则。}
  \label{fig:11.5}
\end{figure}

幸运的是，奇迹确实发生了。我们将在下面看到，\eqref{eq:11.15} 的抵消项虽然只包含三个可调参数，却确实足以抵消该理论中出现的所有无穷大。在本节中，我们将在单圈级显式地演示这种抵消。本章的其余部分将致力于对这些问题的更一般讨论。

\subsection{重整化条件}

在下面的讨论中，我们将只跟踪费曼图的发散部分。然而，记住一组重整化条件是很有用的，它们在原则上也可用于确定抵消项的有限部分。由于抵消项包含三个可调参数，我们需要三个条件。我们可以取这些条件为 \eqref{eq:10.19}（按 \eqref{eq:10.28} 实施），规定 $\sigma$ 场的物理质量 $m$、其场强以及阈散射振幅。然而，在技术上更容易的做法是用对 $\sigma$ 单点振幅（蝌蚪图之和）的约束来替换其中一个条件：
\begin{equation}
  \text{(所有蝌蚪图之和)} = 0 .
  \label{eq:11.16}
\end{equation}
在 QED 中蝌蚪图自动消失，正如我们在式~\eqref{eq:10.5} 中看到的那样。然而，在线性西格马模型中，没有任何对称性禁止非零的单 $\sigma$ 振幅出现。这一振幅会产生 $\sigma$ 的真空期望值，从而，由于 $\phi^N = v + \sigma$，会平移 $\phi$ 的真空期望值。这样的平移是完全可接受的，只要在把抵消项正确地加入振幅计算之后它是有限的。然而，把我们的约定建立为使关系 \eqref{eq:11.16} 在微扰论的各阶都成立，将简化簿记工作。

我们将像式~\eqref{eq:10.19} 那样，把 $\lambda$ 定义为阈散射振幅。于是式~\eqref{eq:11.16} 定义了参数 $\mu$，因此 $\sigma$ 场的质量 $m$ 将与经典方程的结果 $m^2 = 2\mu^2 = 2\lambda v^2$ 相差 $\mathcal{O}(\lambda\mu^2)$ 量级的项。如果我们确实能通过调整三个抵消项来消除理论中的发散，那么这些修正将是有限的，并构成量子场论的一个预言。

概括起来，我们将使用以下重整化条件：
\begin{equation}
  \begin{aligned}
    \text{(蝌蚪图)} &= 0;\qquad
    \frac{d\Sigma}{dp^2}\Big|_{p^2=m^2} = 0;\\
    \text{(两点函数)} &\to \frac{i}{p^2-m^2}
      \text{ 在 } p^2 = m^2 ;\\
    \text{(四点振幅)} &= -i\lambda \text{ 在 } s = 4m^2,\
      t = u = 0 .
  \end{aligned}
  \label{eq:11.17}
\end{equation}
在最后一个条件中，圆圈表示去掉外腿的四点振幅。注意，最后两个条件依赖于 $\sigma$ 粒子的物理质量 $m$。我们现在必须证明，这三个条件足以使线性西格马模型的所有单圈振幅都有限。

\subsection{顶角抵消项}

我们首先通过计算 $4\sigma$ 振幅来确定抵消项 $\delta_\lambda$。树图项来自 $4\sigma$ 顶点，恰好满足 \eqref{eq:11.17}。对这一振幅的单圈贡献是如下图的求和：
\begin{equation}
  \text{(一个 $\sigma$ 圈)} + \text{(一个 $\pi$ 圈)} + \text{(交叉图)} .
  \label{eq:11.18}
\end{equation}
根据 \eqref{eq:11.17}，我们必须调整 $\delta_\lambda$，使这些图的求和在阈处为零。在这个计算中，我们只跟踪紫外发散。这大大简化了分析，因为 \eqref{eq:11.18} 中的大多数图都是有限的。所有圈由三个或更多传播子构成的图都是有限的，因为它们在分母中至少有六次圈动量幂次；例如，
\begin{equation}
  \int \frac{d^4k}{(2\pi)^4}\ \frac{1}{k^2}\frac{1}{k^2}\frac{1}{k^2}
  \quad \text{是有限的} .
\end{equation}
或者，我们可以用下面的方式看出这个图是有限的：每个三点顶点带有一个 $\mu$ 因子，它具有质量量纲。根据第~10.1 节的量纲分析论证，每个这样的因子使图的发散度降低 1。由于 $4\sigma$ 振幅已经有 $D = 0$，任何含三点顶点的图必然是有限的。

我们只剩下 \eqref{eq:11.18} 中的前两个图以及由交叉得到的相关四个图。让我们用维数正规化来评估第一个图：
\begin{equation}
  \begin{aligned}
    \text{(一个 $\sigma$ 圈)} &= \frac{1}{2}(-6i\lambda)^2\ \int
      \frac{d^dk}{(2\pi)^d}\ \frac{i}{k^2-2\mu^2}\ \frac{i}{(k+p)^2-2\mu^2}\\
    &= 18i\lambda^2\ \int_0^1 dx\ \frac{\Gamma(2-d/2)}{(4\pi)^{d/2}}\
      \frac{1}{\Delta^{2-d/2}}\\
    &= 18i\lambda^2\ \frac{\Gamma(2-d/2)}{(4\pi)^2}\ + (\text{有限项}).
  \end{aligned}
  \label{eq:11.19}
\end{equation}
这里 $\Delta$ 是 $p$ 和 $\mu$ 的函数，其精确形式与我们无关。由于我们的目标只是演示发散的抵消，我们将在此处以及本节其余部分忽略有限项。\eqref{eq:11.18} 中的第二个图（内线用 $\pi$ 代替 $\sigma$）与之相同，只是每个顶点因子从 $-6i\lambda$ 变为 $-2i\lambda\delta^{ij}$。（罗马指标 $i, j, \ldots$ 从 1 取到 $N-1$。）因此我们有
\begin{equation}
  \text{(一个 $\pi$ 圈)} = -2i\lambda^2(N-1)\ \frac{\Gamma(2-d/2)}
    {(4\pi)^2}\ + (\text{有限项}) .
  \label{eq:11.20}
\end{equation}
由于这些图每一个的无穷大部分只是一个与动量无关的常数，对应的 $t$ 道和 $u$ 道图的无穷大部分必然相同。因此 $4\sigma$ 顶点的无穷大部分恰好是 \eqref{eq:11.19} 和 \eqref{eq:11.20} 之和的三倍：
\begin{equation}
  \text{($4\sigma$ 顶点的无穷大部分)} \sim 3i\lambda^2(N+8)\
    \frac{\Gamma(2-d/2)}{(4\pi)^2} .
  \label{eq:11.21}
\end{equation}
（在本节中，我们用 $\sim$ 符号表示在忽略有限修正的意义下相等。）应用 \eqref{eq:11.17} 的第三个条件，我们发现抵消项 $\delta_\lambda$ 由下式给出
\begin{equation}
  \delta_\lambda \sim \lambda^2(N+8)\ \frac{\Gamma(2-d/2)}{(4\pi)^2} .
  \label{eq:11.22}
\end{equation}

一旦我们确定了 $\delta_\lambda$ 的值，就固定了另外两个四点振幅的抵消项。这些振幅也都被变得有限了吗？考虑有两个 $\sigma$ 和两个 $\pi$ 的振幅。它从
\begin{equation}
  \text{(一个 $\sigma$ 圈)} + \text{(两个 $\pi$ 圈)} + \text{(一个
    $\sigma\pi$ 圈)} ,
  \label{eq:11.23}
\end{equation}
以及若干带三点顶点的图接受单圈修正，正如前面论证的，这些图显然都是有限的。\eqref{eq:11.23} 中的每个图都包含一个与 \eqref{eq:11.19} 中类似的圈积分，其无穷大部分总是 $-i\Gamma(2-d/2)/(4\pi)^2$。唯一的不同在于顶点和对称因子。例如，\eqref{eq:11.23} 第一个图的无穷大部分是
\begin{equation}
  \sim i\cdot(-6i\lambda)(-2i\lambda\delta^{ij})\cdot
    \frac{1}{2}\ \frac{\Gamma(2-d/2)}{(4\pi)^2}
  = -6i\lambda^2\delta^{ij}\ \frac{\Gamma(2-d/2)}{(4\pi)^2} .
\end{equation}
第二个图稍微复杂一些：
\begin{equation}
  \begin{aligned}
    &\sim \frac{1}{2}(-2i\lambda\delta^{kl})(-2i\lambda(\delta^{ij}\delta^{kl}
      + \delta^{ik}\delta^{jl} + \delta^{il}\delta^{jk}))\cdot
        \frac{\Gamma(2-d/2)}{(4\pi)^2}\\
    &= 2i\lambda^2(N+1)\delta^{ij}\ \frac{\Gamma(2-d/2)}{(4\pi)^2} .
  \end{aligned}
\end{equation}
在第三个图中没有对称因子：
\begin{equation}
  \sim (-2i\lambda\delta^{ik})(-2i\lambda\delta^{kj})\cdot
    \frac{\Gamma(2-d/2)}{(4\pi)^2}
  = -4i\lambda^2\delta^{ij}\ \frac{\Gamma(2-d/2)}{(4\pi)^2} .
\end{equation}
\eqref{eq:11.23} 的第四个图给出相同的表达式，因为它与第三个图相同，只是 $i$ 和 $j$ 互换。因此四个图之和给出 $\sigma\sigma\pi\pi$ 顶点的无穷大部分为
\begin{equation}
  \sim 2i\lambda^2\bigl( 3(N+8) \bigr)\ \frac{\Gamma(2-d/2)}{(4\pi)^2}
  \ \cdots ,
  \label{eq:11.24}
\end{equation}
这个发散项确实被 $\sigma\sigma\pi\pi$ 抵消项所抵消，其中 $\delta_\lambda$ 的取值由 \eqref{eq:11.22} 给出。

剩下的四点振幅有四个外 $\pi$ 场。发散的单圈图是：
\begin{equation}
  \text{(两个 $\pi$ 圈)} + \text{(两个 $\sigma$ 交换图)}
    + \text{(交叉)} .
\end{equation}
这些图都具有同样的熟悉形式。第一个是
\begin{equation}
  \frac{1}{2}(-2i\lambda\delta^{ij})(-2i\lambda\delta^{ij})\cdot
    \frac{\Gamma(2-d/2)}{(4\pi)^2}
  = -2i\lambda^2\delta^{ij}\ \frac{\Gamma(2-d/2)}{(4\pi)^2} ,
  \label{eq:11.25}
\end{equation}
其余的类似。把所有发散图相加并加上抵消项，人们发现 $4\pi$ 振幅被同一个抵消项 $\delta_\lambda$ 变得有限。

\subsection{蝌蚪图与 $\sigma$ 质量}

接下来考虑蝌蚪图和两点函数。蝌蚪（单点）振幅从
\begin{equation}
  \text{(一个 $\sigma$ 圈)} + \text{(一个 $\pi$ 圈)} + \text{(抵消项)} .
  \label{eq:11.30}
\end{equation}
接受贡献。引入 $\pi$ 场的一个小质量 $\xi$ 作为红外正规化子是方便的。第一个图是
\begin{equation}
  \text{(一个 $\sigma$ 圈)} \sim \frac{1}{2}(-6i\lambda v)\ \int
    \frac{d^dk}{(2\pi)^d}\ \frac{i}{k^2-2\mu^2}
  \sim -3i\lambda v\ \frac{\Gamma(1-d/2)}{(4\pi)^{d/2}}\ (2\mu^2)^{d/2-1} .
  \label{eq:11.31}
\end{equation}
第二个图是
\begin{equation}
  \text{(一个 $\pi$ 圈)} \sim \frac{1}{2}(-2i\lambda v\delta^{ii})\ \int
    \frac{d^dk}{(2\pi)^d}\ \frac{i}{k^2-\xi^2}
  \sim -i\lambda v(N-1)\ \frac{\Gamma(1-d/2)}{(4\pi)^{d/2}}\ \xi^{d-2} .
  \label{eq:11.32}
\end{equation}
注意，对于 $d > 2$，该图在 $\xi\to 0$ 极限下消失；然而，它在 $d = 2$ 处有一个极点。尽管有这些奇怪的特征，我们仍可以把 \eqref{eq:11.32} 加到 \eqref{eq:11.31} 上，并要求蝌蚪图被来自图~\ref{fig:11.5} 的抵消项所抵消。这一条件给出
\begin{equation}
  \delta_v + (\text{圈图}) = -\frac{\lambda v}{2}\ \frac{\Gamma(1-d/2)}
    {(4\pi)^{d/2}}\ \Bigl( (2\mu^2)^{1-d/2} + (N-1)\xi^{2-d/2} \Bigr) .
  \label{eq:11.33}
\end{equation}

现在考虑 $2\sigma$ 振幅。单粒子不可约振幅从四个单圈图和一个抵消项接受贡献：
\begin{equation}
  \text{(四个单圈图)} + \text{(抵消项)} .
  \label{eq:11.34}
\end{equation}
把抵消项顶点写为如下形式是方便的
\begin{equation}
  -i\bigl( 2v^2\delta_\lambda \bigr) - i\bigl( \delta_\mu + v^2\delta_\lambda
    \bigr) + ip^2\delta_Z .
  \label{eq:11.35}
\end{equation}
在一般的重整化方案中，$\sigma$ 质量也会被蝌蚪图（及其抵消项）所平移：
\begin{equation}
  \text{(对质量的蝌蚪图修正)} .
  \label{eq:11.36}
\end{equation}
然而，\eqref{eq:11.17} 中的第一个重整化条件迫使这些图精确抵消。这是这一重整化条件特殊简洁性的一个例子。

\eqref{eq:11.34} 的前两个图又是普通四点图的表现形式，只是现在有两条外腿被 $\phi$ 的真空期望值所取代。与前面的计算类比，我们对第一个图得到
\begin{equation}
  \sim 18i\lambda^2v^2\ \frac{\Gamma(2-d/2)}{(4\pi)^2} ,
\end{equation}
对第二个图得到
\begin{equation}
  \sim 2i\lambda^2v^2(N-1)\ \frac{\Gamma(2-d/2)}{(4\pi)^2} .
\end{equation}
利用 \eqref{eq:11.22}，我们看到这两份贡献被 \eqref{eq:11.35} 的第一项所抵消。\eqref{eq:11.34} 的第三个和第四个图包含与 \eqref{eq:11.30} 的蝌蚪图完全相同的积分。关系 \eqref{eq:11.33} 意味着它们被 \eqref{eq:11.35} 中的第二项所抵消。注意，\eqref{eq:11.34} 的任何单圈图中都没有正比于 $p^2$ 的发散项。因此重整化常数 $\delta_Z$ 在单圈级为零，正如在普通 $\phi^4$ 理论中一样。

只剩下一个潜在发散的振幅——$\pi\pi\pi\pi$ 振幅：
\begin{equation}
  \text{(三个单圈图)} + \text{(抵消项)} .
  \label{eq:11.37}
\end{equation}
与 \eqref{eq:11.31} 类比，第一个图是
\begin{equation}
  \sim \frac{1}{2}(-2i\lambda\delta^{ij})\ \frac{\Gamma(1-d/2)}
    {(4\pi)^{d/2}}\ \cdots .
\end{equation}
第二个图与之非常相似。与 \eqref{eq:11.32} 一样，引入一个小 $\pi$ 介子质量作为红外正规化子是有用的。第三个图由下式给出
\begin{equation}
  \sim (-2i\lambda v\delta^{ik})(-2i\lambda v\delta^{kj})\ \int
    \frac{d^dk}{(2\pi)^d}\ \frac{1}{k^2-2\mu^2}\ \frac{1}{(k+p)^2-\xi^2}
  = 4i\lambda^2v^2\delta^{ij}\ \frac{\Gamma(2-d/2)}{(4\pi)^{d/2}}\
    \int_0^1 dx\ \Bigl( \frac{1}{2\mu^2x + \xi^2(1-x) - p^2x(1-x)}
      \Bigr)^{2-d/2} .
  \label{eq:11.38}
\end{equation}
这个表达式的发散部分与 $p$ 无关，所以要检查发散的抵消，只需令 $p = 0$。在 $p = 0$ 处计算完整的振幅（包括有限项）将是有启发性的。把三个圈图与抵消项（其值由 \eqref{eq:11.33} 给出）相加，我们发现 $p = 0$ 处的 $\pi\pi$ 振幅——正是 $\pi$ 场的质量平移 $\delta m_\pi^2$——不仅有限，而且完全为零：
\begin{equation}
  \delta m_\pi^2 \Big|_{p=0} = 0 .
  \label{eq:11.39}
\end{equation}
这个结果非常有吸引力。$p = 0$ 处的 $\pi\pi$ 振幅正是 $\pi$ 场的质量平移 $\delta m_\pi^2$。我们已经知道 $\pi$ 粒子在树图级是无质量的——它们就是戈德斯通定理所要求的 $N-1$ 个无质量玻色子。我们现在已经验证，在线性西格马模型中这些玻色子在单圈级仍然是无质量的；换言之，对线性西格马模型的第一阶量子修正也尊重戈德斯通定理。在本章结尾，我们将给出一个一般论证，说明戈德斯通定理在微扰论的各阶都成立。

\section{有效作用量}

在本章第一节中，我们分析了经典场论中的自发对称性破缺。那一分析是几何性的：我们通过寻找势能面上的最深势阱来找到真空态，并通过表明对称性要求在势阱底部存在一条简并极小值线来证明戈德斯通定理。但这一几何图像在第~11.2 节的单圈计算中丢失了，或者至少被掩盖了。发展一个能让我们即使在量子层面也能运用关于自发对称性破缺的几何论证的形式体系，似乎是值得的。

为更明确地定义我们的目标，考虑确定量子场 $\phi$ 的真空期望值的问题。这个期望值应当作为拉格朗日量参数的函数来确定。在经典层面，计算 $\langle\phi\rangle$ 很容易；只需将势能极小化。然而，正如我们在前一节看到的，这个经典值会被微扰圈修正所改变。事实上，我们看到 $\langle\phi\rangle$ 可能被一个潜在发散的量所平移，而我们需要通过重整化来控制它。

如果在完整的量子场论中存在一个函数，其极小值给出 $\langle\phi\rangle$ 的精确值，那将是极好的。这个函数在微扰论的最低阶将与经典势能一致，但在更高阶会被量子修正所修改。这些修正大概需要重整化来消除无穷大。然而，重整化之后，这个量应当给出 $\langle\phi\rangle$ 与粒子质量和耦合常数之间与我们直接费曼图计算所得相同的那些关系。在本节中，我们将展示一个具有这些性质的函数，称为\emph{有效势}。在第~11.4 节中，我们将解释如何用重整化的质量和耦合常数在微扰论中计算有效势。然后我们将把它用作分析具有隐藏对称性的理论之可重整性的工具。

为识别有效势，考虑第~9.3 节建立的量子场论与统计力学之间的类比。在那节中，我们推导了量子场的关联函数与相关统计系统关联函数之间的对应关系，其中量子涨落被热涨落所取代。在零温下，热力学基态是能量最低的状态，但在非零温度下，我们仍有一个关于优选热力学状态的几何图像：它是使吉布斯自由能取极小值的状态。更明确地说，以磁系统为例，人们通过下式定义亥姆霍兹自由能 $F(H)$
\begin{equation}
  Z(H) = e^{-\beta F(H)} = \int \mathcal{D}s\ \exp\Bigl[ -\beta
    \int d^3x\ \bigl( \mathcal{H}(s) - H s \bigr) \Bigr] .
  \label{eq:11.40}
\end{equation}
系统在给定 $H$ 值处的热力学状态是使吉布斯自由能 $G(M)$ 取极小值的状态，而 $G(M)$ 通过勒让德变换与亥姆霍兹自由能 $F(H)$ 相联系，
\begin{equation}
  G(M) = F(H) - HM ,
\end{equation}
其中 $M = \langle s\rangle$ 是磁化强度。磁化强度是使 $G(M)$ 取极小值的 $s$ 的值。

在量子场论中，我们可以定义这些量的类比物。生成泛函为
\begin{equation}
  Z[J] = e^{-iE[J]} = \int \mathcal{D}\phi\ \exp\Bigl[ i\int d^4x\
    \bigl( \mathcal{L} + J\phi \bigr) \Bigr] ,
  \label{eq:11.43}
\end{equation}
其中 $E[J]$ 是亥姆霍兹自由能的类比物。经典场由下式定义
\begin{equation}
  \phi_{cl}(x) = \langle\phi(x)\rangle_J = \frac{\delta E[J]}
    {\delta J(x)} .
  \label{eq:11.44}
\end{equation}
\emph{有效作用量} $\Gamma[\phi_{cl}]$ 通过勒让德变换定义
\begin{equation}
  \Gamma[\phi_{cl}] = -E[J] - \int d^4x\ J(x)\phi_{cl}(x) ,
  \label{eq:11.47}
\end{equation}
其中通过 \eqref{eq:11.44} 用 $\phi_{cl}$ 消去 $J$。
$\Gamma$ 的泛函导数为
\begin{equation}
  \frac{\delta\Gamma[\phi_{cl}]}{\delta\phi_{cl}(x)} = -J(x) .
  \label{eq:11.48}
\end{equation}
因此，当 $J = 0$ 时，有效作用量被场的真空期望值取极值。对于平移不变的场组态，我们对常数 $\phi_{cl}$ 定义\emph{有效势}
\begin{equation}
  \Gamma[\phi_{cl}] = -\int d^4x\ V_{eff}(\phi_{cl}) ,
  \label{eq:11.50}
\end{equation}
于是 $\phi$ 的真空期望值是使 $V_{eff}$ 取极小值的值：
\begin{equation}
  \frac{\partial V_{eff}}{\partial\phi_{cl}}\Big|_{\phi_{cl} =
    \langle\phi\rangle} = 0 .
  \label{eq:11.51}
\end{equation}
式~\eqref{eq:11.51} 的每个解都是 $J = 0$ 的平移不变状态。式 \eqref{eq:11.47} 意味着在这种情况下 $\Gamma = -E$，因此在 \eqref{eq:11.51} 的一个解处求值的 $V_{eff}(\phi_{cl})$ 正是相应状态的能量密度。

图~\ref{fig:11.6} 说明了函数 $V_{eff}(\phi)$ 的一种可能形状。局部极大值（或者，对于含多个场 $\phi^i$ 的系统，可能是鞍点）是不稳定的组态，不能作为定态实现。图中还包含一个 $V_{eff}$ 的局部极小值，它不是绝对极小值；这是一个\emph{亚稳真空态}，它可以通过量子力学隧穿衰变到真真空。$V_{eff}$ 的绝对极小值是理论中能量最低的状态，因此是真正的、稳定的真空态。具有自发对称性破缺的系统将具有 $V_{eff}$ 的若干个极小值，由于对称性，它们都具有相同的能量。在这些真空中的选择就是自发对称性破缺。

\begin{figure}[t]
  \centering
  \includegraphics[width=0.55\textwidth]{fig11.6.png}
  \caption{标量场理论中有效势的一种可能形式。有效势的极值出现在点 $\phi_{cl} = \phi_1, \phi_2, \phi_3$ 处。真真空态是对应于 $\phi_1$ 的那个。状态 $\phi_2$ 是不稳定的。状态 $\phi_3$ 是亚稳的，但它可以通过量子力学隧穿衰变到 $\phi_1$。}
  \label{fig:11.6}
\end{figure}

\begin{figure}[t]
  \centering
  \includegraphics[width=0.55\textwidth]{fig11.7.png}
  \caption{图~\ref{fig:11.6} 系统的有效势的精确凸形式。}
  \label{fig:11.7}
\end{figure}

在画图~\ref{fig:11.6} 时，我们假设是在对 $\phi$ 的一个固定的常数背景值计算有效势。在某些情况下，对具有给定 $\phi$ 期望值的状态而言，这一状态给出的并非真正的最小能量组态。这种失配可能以下面的方式发生：在一个常数背景场有效势由图~\ref{fig:11.6} 给出的系统中，考虑选择一个介于局部稳定真空态 $\phi_1$ 与 $\phi_3$ 之间的 $\phi_{cl}$ 值：
\begin{equation}
  \phi_{cl} = x\phi_1 + (1-x)\phi_3, \qquad 0 < x < 1 .
  \label{eq:11.52}
\end{equation}
常数背景场的假设给出一个很大的有效势值，如图所示。我们可以通过考虑这样的状态来获得一个能量更低的组态：在这些状态中，存在一些宏观区域 $\langle\phi\rangle = \phi_1$，而另一些区域 $\langle\phi\rangle = \phi_3$，使得 $\langle\phi\rangle$ 在整个系统上的平均值为 $\phi_{cl}$。对于这样的组态，平均真空能由下式给出
\begin{equation}
  V_{eff}(\phi_{cl}) = xV_{eff}(\phi_1) + (1-x)V_{eff}(\phi_3) ,
  \label{eq:11.53}
\end{equation}
如图~\ref{fig:11.7} 所示。内插 \eqref{eq:11.53} 是热力学自由能的麦克斯韦构造在场论中的类比物。一般地，对于满足 \eqref{eq:11.52} 的任意 $\phi_{cl}$、$\phi_1$、$\phi_3$，估计式 \eqref{eq:11.53} 将是有效势的一个上界；我们说有效势是 $\phi_{cl}$ 的\emph{凸函数}。%
\footnote{吉布斯自由能的凸性是统计力学中一个众所周知的精确结果；参见，例如，D.~Ruelle, \emph{Statistical Mechanics} (W.~A.~Benjamin, Reading, Mass., 1969)。}

正如在热力学中一样，计算有效势的直接方案没有考虑相分离的可能性，因此导致如图~\ref{fig:11.6} 所示类型的不稳定和亚稳组态结构。必须手动进行麦克斯韦构造以得到 $V_{eff}(\phi_{cl})$ 的最终形式。幸运的是，$V_{eff}$ 的绝对极小值不受这一细节的影响。

我们现在已经解决了本节开头提出的问题：由式~\eqref{eq:11.47} 和 \eqref{eq:11.50} 定义的有效势给出一个易于可视化的函数，其极小化定义了量子场论的精确真空态，包括量子修正的所有效应。从这些定义出发，如何计算 $V_{eff}$ 并不显然。我们将在下一节中看到，通过直接计算泛函积分来做这件事。

\section{有效作用量的计算}

既然我们已经定义了那个其极小化给出量子场论精确真空态的量，我们就必须学习如何计算它。这可以通过不止一种方式完成。这里我们将使用的最简单方法要求我们有足够的勇气直接从其泛函积分定义来评估完整的有效作用量 $\Gamma$。%
\footnote{此方法归于 R.~Jackiw, \emph{Phys.\ Rev.} \textbf{D9}, 1686 (1974)。}
计算 $\Gamma$ 之后，我们可以通过特殊化到 $\phi_{cl}$ 的常数值来得到 $V_{eff}$。

我们的计划是为生成泛函 $Z$ 找到一个微扰展开，从其泛函积分定义 \eqref{eq:11.43} 出发。然后取对数得到能量泛函 $E$，最后按照式~\eqref{eq:11.47} 做勒让德变换得到 $\Gamma$。我们将使用重整化微扰论，所以像式~\eqref{eq:10.18} 中那样把拉格朗日量分成依赖于重整化参数的部分和包含抵消项的部分是方便的：
\begin{equation}
  \mathcal{L} = \mathcal{L}_1 + \delta\mathcal{L} .
  \label{eq:11.54}
\end{equation}
我们希望把 $\Gamma$ 计算为 $\phi_{cl}$ 的函数。但泛函 $Z[J]$ 通过其对 $J$ 的依赖而依赖于 $\phi_{cl}$。因此，我们必须至少隐含地找到 $J(x)$ 与 $\phi_{cl}(x)$ 之间的关系。在微扰论的最低阶，这一关系正是经典场方程：
\begin{equation}
  0 = \phi_{cl} - \frac{\delta}{\delta J} \cdots
  \qquad\text{(到最低阶)} .
\end{equation}
让我们把 $J_1(x)$ 定义为当 $\mathcal{L} = \mathcal{L}_1$ 时精确满足此方程的任意函数：
\begin{equation}
  0 = \frac{\delta\mathcal{L}_1}{\delta\phi}\Big|_{\phi=\phi_{cl}}
    + J_1(x) \qquad \text{(精确地)} .
  \label{eq:11.55}
\end{equation}
我们将把 $J$ 与 $J_1$ 之间的差看作一个抵消项，类似于 $\delta\mathcal{L}$，所以我们写
\begin{equation}
  J(x) = J_1(x) + \delta J(x) ,
  \label{eq:11.56}
\end{equation}
其中 $\delta J$ 由 $\phi_{cl}$ 的原始定义 \eqref{eq:11.44}，即 $\langle\phi(x)\rangle_{J=0} = \phi_{cl}(x)$，在微扰论中逐阶确定。

利用这种记号，我们把式~\eqref{eq:11.43} 重写为
\begin{equation}
  Z[J] = \int \mathcal{D}\phi\ \exp\Bigl[ i\int d^4x\ \bigl(
    \mathcal{L}_1 + J_1\phi \bigr)\ \Bigr]\ \exp\Bigl[ i\int d^4x\ \bigl(
      \delta\mathcal{L} + \delta J\phi \bigr) \Bigr] .
  \label{eq:11.57}
\end{equation}
第二个指数中包含抵消项；暂时把它放在一边。在第一个指数中，通过替换 $\phi(x) = \phi_{cl}(x) + \eta(x)$ 把指数在 $\phi_{cl}$ 附近展开。这个指数具有如下形式
\begin{equation}
  \begin{aligned}
    \int d^4x\ \bigl( \mathcal{L}_1 + J_1\phi \bigr)
    &= \int d^4x\ \bigl( \mathcal{L}_1[\phi_{cl}] + J_1\phi_{cl} \bigr)
      + \int d^4x\ \eta(x)\Bigl( \frac{\delta\mathcal{L}_1}{\delta\phi}
        + J_1 \Bigr)\\
    &\quad + \frac{1}{2}\int d^4x\,d^4y\ \eta(x)\eta(y)\
      \frac{\delta^2\mathcal{L}_1}{\delta\phi(x)\delta\phi(y)}
      + \cdots ,
  \end{aligned}
  \label{eq:11.58}
\end{equation}
其中 $\mathcal{L}_1$ 的各种泛函导数都在 $\phi_{cl}(x)$ 处取值。注意，线性于 $\eta$ 的项通过使用式~\eqref{eq:11.55} 而消失。因此对 $\eta$ 的积分是一个高斯积分，三次及更高次项给出微扰修正。

我们将按照第~9.2 节的法则描述对这一积分的正规求值。这一求值中的要素将是式~\eqref{eq:11.58} 的系数，即 $\mathcal{L}_1$ 的各阶泛函导数。暂时请接受它们给出定义良好的算符。在给出 $\Gamma[\phi_{cl}]$ 的一般表达式之后，我们将在一个标量场论例子中显式地进行这一计算。我们将在这个例子中看到，这些形式算符对应于费曼图微扰论中熟悉的表达式。

那么，让我们考虑使用展开式 \eqref{eq:11.58} 来对 $\eta(x)$ 积分。只保留到 $\eta$ 的二次项，并且仍然忽略抵消项，我们得到一个纯高斯积分，它可以用泛函行列式来求值：
\begin{equation}
  \int \mathcal{D}\eta\ \exp\Bigl[ \frac{i}{2}\int d^4x\,d^4y\
    \eta(x)\,\frac{\delta^2\mathcal{L}_1}{\delta\phi(x)\delta\phi(y)}\,
    \eta(y) \Bigr] = \Bigl[ \det\Bigl( -i\frac{\delta^2\mathcal{L}_1}
      {\delta\phi(x)\delta\phi(y)} \Bigr) \Bigr]^{-1/2} .
  \label{eq:11.59}
\end{equation}
这个泛函行列式将给出有效作用量的最低阶量子修正，并且对许多目的而言，没有必要在展开式 \eqref{eq:11.58} 中继续推进。稍后我们将看到，如果我们确实包含 $\eta$ 的三次及更高次项，它们会产生泛函积分 \eqref{eq:11.57} 的费曼图展开，其中传播子为算符的逆
\begin{equation}
  \Bigl( -i\frac{\delta^2\mathcal{L}_1}{\delta\phi(x)\delta\phi(y)}
    \Bigr)^{-1} ,
  \label{eq:11.60}
\end{equation}
顶点为 $\mathcal{L}_1$ 的三阶及更高阶泛函导数。

最后，让我们放回式~\eqref{eq:11.57} 中第二个指数的影响，即抵消项拉格朗日量。把这个项在 $\phi = \phi_{cl}$ 附近展开是有用的，把它写为
\begin{equation}
  \bigl( \delta\mathcal{L}[\phi_{cl}] + \delta J\phi_{cl} \bigr)
  + \bigl( \delta\mathcal{L}[\phi_{cl}+\eta] - \delta\mathcal{L}[\phi_{cl}]
    + \delta J\eta \bigr) .
  \label{eq:11.61}
\end{equation}
\eqref{eq:11.61} 的第二项可以展开为 $\eta$ 的泰勒级数；逐次项给出抵消项顶点，它们可以被包含在前面提到的费曼图中。第一项相对于对 $\eta$ 的泛函积分是常数，因此会在式~\eqref{eq:11.59} 的指数中给出附加项。

把积分 \eqref{eq:11.59} 与来自高阶顶点和抵消项的贡献相结合，可以得到泛函积分 \eqref{eq:11.57} 的完整表达式。我们将在下面的例子中看到，表示高阶项的费曼图可以被安排成给出连通图之和的指数。于是得到 $E[J]$ 的如下表达式：
\begin{equation}
  \begin{aligned}
    -iE[J] &= i\int d^4x\ \bigl( \mathcal{L}_1[\phi_{cl}] + J\phi_{cl} \bigr)
      - \frac{i}{2}\log\det\Bigl[ -i\frac{\delta^2\mathcal{L}_1}
        {\delta\phi(x)\delta\phi(y)} \Bigr]\\
    &\quad + (\text{连通图}) + i\int d^4x\ \bigl(
      \delta\mathcal{L}[\phi_{cl}] + \delta J\phi_{cl} \bigr) .
  \end{aligned}
  \label{eq:11.62}
\end{equation}
从这个方程可以直接得到 $\Gamma$：利用 $\delta E/\delta\phi_{cl} + \delta J = J$ 和勒让德变换 \eqref{eq:11.47}，我们得到
\begin{equation}
  \Gamma[\phi_{cl}] = \int d^4x\ \mathcal{L}_1[\phi_{cl}]
    + \frac{i}{2}\log\det\Bigl[ -i\frac{\delta^2\mathcal{L}_1}
      {\delta\phi(x)\delta\phi(y)} \Bigr]
    - i\cdot(\text{连通图}) + \int d^4x\ \delta\mathcal{L}[\phi_{cl}].
  \label{eq:11.63}
\end{equation}
注意，没有任何项显式地依赖于 $J$ 留下来；因此，$\Gamma$ 被表示为 $\phi_{cl}$ 的函数，理当如此。对 $\Gamma[\phi_{cl}]$ 有贡献的费曼图没有外线，最简单的图恰好是两圈图。$\Gamma$ 的最低阶量子修正由泛函行列式给出，而这一项就是本书将要用到的全部。

\eqref{eq:11.63} 的最后一项提供了一组抵消项，可用于满足对 $\Gamma$ 的重整化条件，并在此过程中抵消在泛函行列式和图的计算中出现的发散。我们将在下面的例子中精确地展示这种抵消如何工作。重整化条件将确定 $\delta\mathcal{L}$ 中所有的抵消项。然而，我们构造的形式体系包含一个新的抵消项 $\delta J$。这个系数由以下特殊判据确定：在式~\eqref{eq:11.55} 中，我们这样建立分析，使得在领头阶 $\langle\phi\rangle = \phi_{cl}$。然而，这一关系有可能在更高阶失效：量 $\langle\phi\rangle$ 可能从费曼图接受附加贡献，从而把它从值 $\phi_{cl}$ 平移开。如果有非零的蝌蚪图对 $\langle\eta\rangle$ 有贡献，就会发生这种情况。但这个振幅也从 \eqref{eq:11.61} 中的抵消项 $(\delta J\eta)$ 接受贡献。因此，我们可以通过调整 $\delta J$ 使其满足图方程
\begin{equation}
  \delta J\ \bigl( \text{蝌蚪图} \bigr)
  = -\bigl( \text{所有蝌蚪图之和} \bigr) = 0 .
  \label{eq:11.64}
\end{equation}
来保持 $\langle\eta\rangle = 0$，并在此过程中把 $\delta J$ 确定到任意阶。在实践中，我们将通过简单地忽略任何单粒子不可约单点图来满足这一条件，因为任何这样的图都会被 $\delta J$ 的调整所抵消。这些蝌蚪图的消除，在第~11.2 节中我们需要一些努力来安排，在这里则作为形式体系的自然组成部分而被内置。

\subsection{线性西格马模型中的有效作用量}

在式~\eqref{eq:11.63} 中，我们给出了 $\Gamma[\phi_{cl}]$ 的一个完整但并非十分透明的求值。现在让我们通过在线性西格马模型中计算 $\Gamma[\phi_{cl}]$ 来阐明这个方程的意义，并把它派上一些好用场。我们将看到，第~11.2 节中我们通过蛮力微扰论得到的结果从式~\eqref{eq:11.63} 中要自然得多地涌现出来。

我们再次从拉格朗日量 \eqref{eq:11.5} 出发：
\begin{equation}
  \mathcal{L} = \tfrac{1}{2}(\partial_\mu\phi^i)^2 + \tfrac{1}{2}\mu^2
    (\phi^i)^2 - \frac{\lambda}{4}\bigl[ (\phi^i)^2 \bigr]^2 .
  \label{eq:11.65}
\end{equation}
在经典场附近展开：$\phi^i = \phi^i_{cl} + \eta^i$。因为我们期望找到平移不变的真空态，我们将特殊化到常数经典场的情形。这将简化下面计算的某些要素。特别是，根据式~\eqref{eq:11.50}，最终结果将正比于泛函积分的四维体积 $(VT)$。当这一依赖被分解出去后，我们将得到有效势的一个定义良好的内含量表达式。无论如何，在这一简化之后，\eqref{eq:11.65} 具有如下形式
\begin{equation}
  \mathcal{L} = \frac{1}{2}\bigl( \mu^2(\phi^i_{cl})^2
    - \frac{\lambda}{2}(\phi^i_{cl})^2)^2 \bigr)
    + \frac{1}{2}\eta^i\bigl( -\partial^2\delta^{ij} + (\lambda(\phi_{cl})^2
      - \mu^2)\delta^{ij} + 2\lambda\phi^i_{cl}\phi^j_{cl} \bigr)\eta^j
    + \cdots .
  \label{eq:11.66}
\end{equation}
二次项定义了我们必须求其行列式的算符。对于常数经典场的情形，我们可以选择坐标系使 $\phi_{cl}$ 指向第 $N$ 个方向，如 \eqref{eq:11.7}：$\phi^i_{cl} = (0, \ldots, 0, \phi_{cl})$。于是算符是对角的，其特征值为
\begin{equation}
  m^2_\sigma = \mu^2 - 3\lambda\phi_{cl}^2
  \qquad\text{和}\qquad
  m^2_\pi = \mu^2 - \lambda\phi_{cl}^2 ,
  \label{eq:11.69}
\end{equation}
分别对应 $\sigma$ 场和 $\pi$ 场。注意，当 $\phi_{cl}^2 = \mu^2/\lambda$（经典真空期望值）时，$\pi$ 特征值为零。这就是有效势层面的戈德斯通定理：质量矩阵恰好在该极小值处有一个零特征值。

\eqref{eq:11.63} 中的泛函行列式于是为
\begin{equation}
  \frac{i}{2}\log\det\Bigl[ -i\frac{\delta^2\mathcal{L}_1}
    {\delta\phi(x)\delta\phi(y)} \Bigr]
  = \frac{i}{2}\log\det\bigl[ -i(\partial^2 + M^2) \bigr],
  \label{eq:11.71}
\end{equation}
其中 $M^2$ 是特征值由 \eqref{eq:11.69} 给出的质量矩阵。利用恒等式 $\log\det M = \mathrm{tr}\log M$ 并对所得高斯泛函积分求值，我们得到
\begin{equation}
  V^{(1)}_{eff}(\phi_{cl}) = \frac{1}{2(4\pi)^{d/2}}\ \Gamma(-d/2)\
    \Bigl[ (m_\sigma^2)^{d/2} + (N-1)(m_\pi^2)^{d/2} \Bigr] .
  \label{eq:11.73}
\end{equation}
像式~\eqref{eq:11.15} 中那样加上抵消项拉格朗日量，我们发现
\begin{equation}
  \begin{aligned}
    V_{eff}(\phi_{cl}) &= -\frac{1}{2}\mu^2\phi_{cl}^2 + \frac{\lambda}{4}
      \phi_{cl}^4 + \frac{1}{2(4\pi)^{d/2}}\ \Gamma(-d/2)\ \Bigl[
        \bigl( \mu^2 - 3\lambda\phi_{cl}^2 \bigr)^{d/2}\\
    &\quad + (N-1)\bigl( \mu^2 - \lambda\phi_{cl}^2 \bigr)^{d/2} \Bigr]
      + 2\delta_\mu\phi_{cl}^2 + 4\delta_\lambda\phi_{cl}^4 .
  \end{aligned}
  \label{eq:11.74}
\end{equation}
这里我们把 $(\phi_{cl})^2$ 写为 $(\phi^i_{cl})^2$ 的简写。由于这个结果的第二行是领头辐射修正，我们或许可以预期该结果具有单圈费曼图的结构。确实，我们看到这个表达式包含伽马函数和紫外发散，类似于我们在第~11.2 节单圈计算中发现的那些。我们将在下面证明，这一项实际上具有与我们在第~11.2 节发现的完全相同的紫外发散。这些发散将被式~\eqref{eq:11.74} 最后一行的抵消项所减去。

\begin{figure}[t]
  \centering
  \includegraphics[width=0.5\textwidth]{fig11.8.png}
  \caption{对 $O(N)$ 线性西格马模型有效势求值有贡献的费曼图：（a）一个被 \eqref{eq:11.64} 移除的图；（b）第一批非零的图修正。}
  \label{fig:11.8}
\end{figure}

由于式~\eqref{eq:11.63} 中行列式的计算给出单圈修正的效应，我们或许可以预期对式~\eqref{eq:11.63} 有贡献的费曼图从两圈阶开始。对于 $O(N)$ 西格马模型的情形，我们可以明确地看到这一点。式~\eqref{eq:11.60} 下面描述的微扰展开涉及传播子，它是式~(11.67) 的逆：
\begin{equation}
  \langle \eta^i(k)\eta^j(-k) \rangle = \frac{i\delta^{ij}}{k^2 - m^2} ,
  \label{eq:11.75}
\end{equation}
其中 $m^2$ 由 \eqref{eq:11.69} 给出。顶点由拉格朗日量展开中 $\eta^3$ 和 $\eta^4$ 阶的项给出。综合这些要素，我们发现对真空能有贡献的领头费曼图具有图~\ref{fig:11.8} 所示的形状。图~\ref{fig:11.8}(a) 中的图实际上被抵消项 $\delta J$ 的效应所抵消，如式~\eqref{eq:11.64} 所示。因此对有效势的领头图贡献来自图~\ref{fig:11.8}(b) 的两圈图。

结果 \eqref{eq:11.74} 显然具有 $O(N)$ 对称性。从我们在第~11.2 节开头提出的问题出发，我们或许曾担心当我们围绕一个具有自发对称性破缺的状态计算辐射修正时，这一性质会被破坏。但 $V_{eff}(\phi_{cl})$ 是我们极小化以寻找真空态的函数，因此即使最低能量真空是不对称的，它也应该恰当地保持对称。在我们这里构造的形式体系中，没有必要担心。公式 \eqref{eq:11.63} 在拉格朗日量的原始 $O(N)$ 对称性下逐项明显不变。因此我们必然得到 $\Gamma_{eff}(\phi_{cl})$ 的一个 $O(N)$ 对称结果。

在继续精确确定 $\delta_\mu$ 和 $\delta_\lambda$ 之前，我们不妨先检查式~\eqref{eq:11.74} 中的抵消项是否足以使 $\Gamma[\phi_{cl}]$ 的表达式有限。因子 $\Gamma(-d/2)$ 在 $d = 0, 2, 4$ 处有极点。$d = 0$ 处的极点是一个常数，与 $\phi_{cl}$ 无关，因此没有物理意义。$d = 2$ 处的极点是 $\phi_{cl}$ 的一个偶二次多项式。$d = 4$ 处的极点是 $\phi_{cl}$ 的一个偶四次多项式。因此，如果令
\begin{equation}
  \delta_\mu = -\lambda\mu^2(N+2)\ \frac{\Gamma(-d/2)}{(4\pi)^{d/2}}
    + \text{有限} .
\end{equation}
式~\eqref{eq:11.74} 在 $d\to 2$ 极限下成为有限表达式。如果令
\begin{equation}
  \delta_\mu = -\lambda\mu^2(N+2)\ \frac{\Gamma(2-d/2)}{(4\pi)^2}
    + \text{有限};\qquad
  \delta_\lambda = \lambda^2(N+8)\ \frac{\Gamma(2-d/2)}{(4\pi)^2}
    + \text{有限} .
  \label{eq:11.76}
\end{equation}
表达式在 $d\to 4$ 时有限。这些表达式分别与第~11.2 节的早期结果式~\eqref{eq:11.33} 和 \eqref{eq:11.22} 在 $d\to 2$ 和 $d\to 4$ 极限下一致。

$\delta_\lambda$ 和 $\delta_\mu$ 的有限部分取决于所施加的重整化条件的具体形式。例如，在第~11.2 节中，我们施加了条件 \eqref{eq:11.16}，即 $\phi$ 的真空期望值等于 $\mu/\sqrt{\lambda}$，以及 \eqref{eq:11.17} 中关于 $\sigma$ 的散射振幅和场强的附加条件。条件 \eqref{eq:11.16} 很容易用有效势表示为
\begin{equation}
  \frac{dV_{eff}}{d\phi_{cl}}\Big|_{\phi_{cl} = \mu/\sqrt{\lambda}} = 0 .
\end{equation}
利用 $\Gamma$ 的导数与单粒子不可约振幅之间的联系，我们可以把另外两个条件写成 $\Gamma[\phi_{cl}]$ 泛函导数的动量空间傅里叶变换。通过这种方式，在原则上可以重构第~11.2 节所使用的特定重整化方案。

然而，如果我们想要可视化由量子修正引起的对最低阶结果的修正，我们可以应用一个更容易实施的重整化方案。这样一种方案，称为\emph{最小减除}（MS），就是简单地去除潜在发散量中的 $(1/\epsilon)$ 极点（其中 $\epsilon = 4-d$）。不过，通常这些 $(1/\epsilon)$ 极点伴随有涉及 $\gamma$ 和 $\log(4\pi)$ 的项。把这些项也一并减去是方便的，也并非更随意。在这种称为\emph{修正的最小减除}或 $\overline{MS}$（``em-ess-bar''）的处方中，人们替换
\begin{equation}
  \frac{1}{(4\pi)^{d/2}(m^2)^{2-d/2}}\ \Gamma(2-d/2)
  \to \frac{1}{(4\pi)^2}\ \Bigl( \frac{2}{\epsilon} - \gamma_E
    + \log 4\pi - \log m^2 \Bigr)
  \to \frac{1}{(4\pi)^2}\ \Bigl( -\log\Bigl( \frac{m^2}{M^2} \Bigr)
    + \cdots \Bigr),
  \label{eq:11.77}
\end{equation}
其中 $M$ 是我们引入的一个任意质量参数，以使最终方程量纲正确。你应该把 $M$ 想成是参数化了一序列可能的重整化条件。$\overline{MS}$ 重整化方案通常使单圈修正具有特别简单的形式。这种简单性的代价是，通常需要一些努力才能用 $\overline{MS}$ 表达式的参数来表示物理可测量量。

要把 $\overline{MS}$ 重整化处方应用于 \eqref{eq:11.74}，我们需要把这个方程中的发散项按 $\epsilon$ 的幂次展开。作为例子，考虑表达式 \eqref{eq:11.73} 的 $\overline{MS}$ 正规化：
\begin{equation}
  \begin{aligned}
    \frac{1}{2(4\pi)^{d/2}}\ \Gamma(-d/2)\,m^{d/2}
    &= \frac{1}{(4\pi)^{d/2}}\ \frac{\Gamma(2-d/2)}{(d/2)(d/2-1)}\ m^2\
      (m^2)^{d/2-2}\\
    &= \frac{m^2}{2(4\pi)^2}\ \Bigl( -\frac{2}{\epsilon} - \log\Bigl(
      \frac{m^2}{M^2} \Bigr) + \frac{3}{2} \Bigr) .
  \end{aligned}
  \label{eq:11.78}
\end{equation}
以这种方式修改我们的结果 \eqref{eq:11.74}，我们发现
\begin{equation}
  \begin{aligned}
    V_{eff} &= -\frac{1}{2}\mu^2\phi_{cl}^2 + \frac{\lambda}{4}\phi_{cl}^4
      + \frac{1}{64\pi^2}\ \Bigl\{ (N-1)\bigl( \mu^2 - \lambda\phi_{cl}^2
        \bigr)^2\Bigl[ \log\Bigl( \frac{\mu^2 - \lambda\phi_{cl}^2}{M^2}
          \Bigr) - \frac{3}{2} \Bigr]\\
    &\quad + \bigl( \mu^2 - 3\lambda\phi_{cl}^2 \bigr)^2\Bigl[ \log\Bigl(
      \frac{\mu^2 - 3\lambda\phi_{cl}^2}{M^2} \Bigr) - \frac{3}{2}
      \Bigr] \Bigr\} .
  \end{aligned}
  \label{eq:11.79}
\end{equation}
有效势从而被修改为在 $\phi_{cl}$ 大值处稍陡、在较小值处更负，如图~\ref{fig:11.9} 所示。对每一组 $\mu$、$\lambda$ 和 $M$ 的值，我们都可以通过对 $\phi_{cl}$ 极小化 $V_{eff}(\phi)$ 来确定优选真空态。当对数的自变量变为负值时，对 $V_{eff}$ 的修正未定义，但幸运的是 $V_{eff}$ 的极小值出现在这一区域之外，如图所示。

\begin{figure}[t]
  \centering
  \includegraphics[width=0.5\textwidth]{fig11.9.png}
  \caption{$\phi^4$ 理论（$N=1$）的有效势，其中量子修正按式~\eqref{eq:11.79} 包含。较浅的曲线显示经典势能，以供比较。}
  \label{fig:11.9}
\end{figure}

在继续之前，我们想对这个有效势表达式提出两个问题。我们将提出的问题在量子场论计算中普遍出现，但表达式 \eqref{eq:11.79} 为这些困难提供了一个具体的例证。我们在接下来两章中的大部分讨论将致力于建立一个能够回答这些问题的形式体系。

首先，令人不安的是，虽然我们的经典拉格朗日量只包含两个参数 $\mu$ 和 $\lambda$，但结果 \eqref{eq:11.79} 依赖于三个参数，其中一个是任意质量标度 $M$。对这种抱怨可以给出一个表面上的回答，如下：考虑把 $M^2$ 的值改为 $M^2 + \delta M^2$ 所导致的 $V_{eff}(\phi_{cl})$ 的变化。从 \eqref{eq:11.79} 的显式形式可以看出，这一变化通过移动 $\mu$ 和 $\lambda$ 的值而被完全补偿，即按
\begin{equation}
  \mu^2 \to \mu^2\Bigl( 1 + \frac{\lambda(N+2)}{16\pi^2}\ \log\Bigl(
    \frac{M^2}{M_0^2} \Bigr) \Bigr),\qquad
  \lambda \to \lambda\Bigl( 1 + \frac{\lambda(N+8)}{16\pi^2}\ \log\Bigl(
    \frac{M^2}{M_0^2} \Bigr) \Bigr) .
  \label{eq:11.80}
\end{equation}
因此，$M^2$ 的改变完全等价于参数 $\mu$ 和 $\lambda$ 的改变。然而，为什么会如此，或者这一事实如何帮助我们理解公式对 $M^2$ 的依赖，这一点并不清楚。

第二个问题源于这样一个事实：式~\eqref{eq:11.79} 中的单圈修正包含一个对数，它可以变得足够大以补偿小耦合常数 $\lambda$。这一问题在 $\mu^2\to 0$ 极限下尤为清楚；此时式~\eqref{eq:11.79} 取如下形式
\begin{equation}
  \begin{aligned}
    V_{eff} &= \frac{\lambda}{4}\phi_{cl}^4 + \frac{1}{64\pi^2}\ \Bigl[
      \lambda^2\phi_{cl}^4\Bigl( (N+8)\Bigl( \log\Bigl(
        \frac{\lambda\phi_{cl}^2}{M^2} \Bigr) - 1 \Bigr) + 9\log 3 \Bigr)
      \Bigr]\\
    &= \frac{\lambda}{4}\phi_{cl}^4\ \Bigl[ 1 + \frac{\lambda}{16\pi^2}\ \Bigl(
      (N+8)\Bigl( \log\Bigl( \frac{\lambda\phi_{cl}^2}{M^2} \Bigr)
        - 1 \Bigr) + 9\log 3 \Bigr) \Bigr] .
  \end{aligned}
  \label{eq:11.81}
\end{equation}
这个势的极小值在哪里？如果我们按这个表达式的表面值来取，我们会发现当 $\phi_{cl}$ 达到一个非常小的、量级为
\begin{equation}
  \phi_{cl} \sim \frac{M}{\sqrt{\lambda}}\ \exp\Bigl[ -\frac{16\pi^2}
    {(N+8)\lambda} \Bigr] ,
\end{equation}
的值时 $V_{eff}(\phi_{cl})$ 经过零点，并且在这一点的附近，在 $\phi_{cl}$ 的非零值处达到一个极小值。但这个零点是通过领头项与量子修正的相消而发生的。换言之，在我们能够处理 $\mu^2 = 0$ 时 $V_{eff}(\phi_{cl})$ 是否有对称性破缺极小值的问题之前，微扰论已经完全失效。看来我们目前的工具完全不足以解决这种情况。

虽然这一点远非显然，但事实证明这两个问题是相互关联的。我们在第~12 章的主要结果之一将是对式~\eqref{eq:11.80} 所显示的 $M^2$、$\lambda$ 和 $\mu^2$ 相互关系的解释。然后，在第~13 章中，我们将利用从这一分析中获得的洞见来完全解决大对数出现的第二个问题。然而，在开始那项研究之前，在具有自发对称性破缺的理论之重整化的更形式化方面，还有一些问题我们尚未讨论。

\section{作为生成泛函的有效作用量}

既然我们已经定义了有效作用量，并为某一个特定理论计算了它，让我们回到理解具有隐藏对称性的理论之重整化的目标上来。在第~11.6 节中，我们将把有效作用量用作实现这一目标的工具。然而，首先我们必须更详细地研究有效作用量与费曼图之间的关系。

我们在第~9.2 节看到，$Z[J]$ 对 $J(x)$ 的泛函导数产生标量场的关联函数（参见，例如，式~\eqref{eq:9.35}）。换言之，$Z[J]$ 是关联函数的生成泛函。我们现在目标是证明 $\Gamma[\phi_{cl}]$ 也是一个这样的生成泛函；具体来说，它是\emph{单粒子不可约}（1PI）关联函数的生成泛函。由于 1PI 关联函数在重整化理论中占有突出地位，这一结果将是下一节重整化讨论的核心。

首先，让我们考虑的不是 $\Gamma[\phi_{cl}]$ 的泛函导数，而是 $E[J] = -i\log Z[J]$ 的泛函导数。一阶导数，由式 \eqref{eq:11.44} 给出，正是 $\langle\phi(x)\rangle$。二阶导数为
\begin{equation}
  \frac{\delta^2 E[J]}{\delta J(x)\delta J(y)}
  = -i\Bigl( \langle\phi(x)\phi(y)\rangle
    - \langle\phi(x)\rangle\langle\phi(y)\rangle \Bigr) ,
  \label{eq:11.82}
\end{equation}
这就是连通两点函数（除了因子 $-i$）。$E[J]$ 的更高阶导数给出连通的 $n$ 点函数：
\begin{equation}
  \frac{\delta^n E[J]}{\delta J(x_1)\cdots\delta J(x_n)} = -i\,
    \langle\phi(x_1)\cdots\phi(x_n)\rangle_{\mathrm{connected}} .
  \label{eq:11.83}
\end{equation}
因此，$E[J]$ 是连通关联函数的生成泛函。

现在，有效作用量通过勒让德变换 \eqref{eq:11.47} 与 $E[J]$ 相联系。$\Gamma$ 的一阶泛函导数为
\begin{equation}
  \frac{\delta\Gamma[\phi_{cl}]}{\delta\phi_{cl}(x)} = -J(x) .
  \label{eq:11.84}
\end{equation}
二阶导数与连通两点函数的逆相联系。对 $\phi_{cl}$ 微分 \eqref{eq:11.44}，我们得到
\begin{equation}
  \delta(x-y) = \int d^4z\ \frac{\delta^2 E[J]}{\delta J(x)\delta J(z)}\
    \frac{\delta J(z)}{\delta\phi_{cl}(y)}
  = \int d^4z\ \frac{\delta^2 E[J]}{\delta J(x)\delta J(z)}\
    \Bigl( -\frac{\delta^2\Gamma}{\delta\phi_{cl}(z)\delta\phi_{cl}(y)}
      \Bigr) .
  \label{eq:11.90}
\end{equation}
利用 \eqref{eq:11.82}，这说明 $\Gamma$ 的二阶导数是连通两点函数的逆（带有适当的 $i$ 因子）。用图形的方式，我们可以写
\begin{equation}
  \frac{\delta^2\Gamma[\phi_{cl}]}{\delta\phi_{cl}(x)\delta\phi_{cl}(y)}
  = iD^{-1}(x-y) ,
  \label{eq:11.91}
\end{equation}
其中 $D(x-y)$ 是完整传播子（连通两点函数）。因此，有效作用量的二阶导数是\emph{逆传播子}，即去掉外传播子后的完整两点函数。

对于三点函数，我们可以微分 \eqref{eq:11.91} 并利用链式法则得到
\begin{equation}
  \frac{\delta^3 E[J]}{\delta J(x)\delta J(y)\delta J(z)}
  = i\int d^4u\,d^4v\,d^4w\ D(x-u)D(y-v)D(z-w)\
    \frac{\delta^3\Gamma[\phi_{cl}]}
      {\delta\phi_{cl}(u)\delta\phi_{cl}(v)\delta\phi_{cl}(w)} .
  \label{eq:11.95}
\end{equation}
这一关系用图形的方式表达更为清晰。左边是连通三点函数。如果我们按 \eqref{eq:11.95} 所示的提取出精确传播子，它就如下分解：$i\Gamma[\phi_{cl}]$ 的三阶导数正是去掉全部三个完整传播子后的连通关联函数，即单粒子不可约三点函数：
\begin{equation}
  \frac{\delta^3\Gamma[\phi_{cl}]}{\delta\phi_{cl}(x)\delta\phi_{cl}(y)
    \delta\phi_{cl}(z)} = \text{(1PI 三点函数)} .
\end{equation}
通过类似（即便越来越复杂）的操作，人们可以对 $\Gamma$ 的每一阶导数推导出同样的关系。例如，微分式~\eqref{eq:11.95}，我们最终发现（使用矩阵记号，重复指标隐式地对之积分）$i\Gamma$ 的四阶导数是单粒子不可约四点函数。一般关系（对于 $n \ge 3$）为
\begin{equation}
  \frac{\delta^n\Gamma[\phi_{cl}]}{\delta\phi_{cl}(x_1)\cdots
    \delta\phi_{cl}(x_n)} = \text{(1PI } n\text{ 点函数)} .
  \label{eq:11.96}
\end{equation}
换言之，有效作用量是单粒子不可约关联函数的生成泛函。

这一结论意味着 $\Gamma$ 包含量子场论的完整物理预言集合。让我们回顾这些信息是如何编码的。场论的真空态被识别为有效势的极小值。极小值的位置决定拉格朗日量的对称性是被保持还是被自发破缺。$\Gamma$ 的二阶导数是逆传播子。传播子的极点，或逆传播子的零点，给出粒子质量的值。因此粒子质量 $m^2$ 被确定为求解以下方程的 $p^2$ 的值
\begin{equation}
  \int d^4x\ e^{ip\cdot(x-y)}\ \frac{\delta^2\Gamma}{\delta\phi(x)
    \delta\phi(y)} = 0 .
  \label{eq:11.97}
\end{equation}
$\Gamma$ 的更高阶导数是单粒子不可约振幅。它们可以用完整传播子连接起来，组合在一起构造四点及更高点的连通振幅，从而给出 $S$ 矩阵元。因此，从对 $\Gamma$ 的认识出发，我们可以重构量子场论的定性行为、其对称性破缺的模式，以及其粒子及其相互作用的定量细节。

\section{重整化与对称性：一般分析}

在我们对量子场论发散的讨论中（尤其是式~\eqref{eq:10.4} 下面那一段），我们注意到费曼积分的基本发散与单粒子不可约图相联系。因此我们或许可以预期，有效作用量在讨论量子场论的可重整性时，尤其是那些具有自发对称性破缺的理论，将是一个有用的对象。在本节中，我们将正好以这种方式利用有效作用量。

在第~11.4 节中，我们在一个特定例子中看到，计算有效作用量的形式体系提供了消除紫外发散所需的抵消项，至少在单圈级如此。这些抵消项恰好是原始拉格朗日量的那些。我们现在将通过把第~10.1 节的幂次计数论证直接应用于有效作用量的计算，来论证这一组抵消项总是充分的——对任意阶以及任何可重整化场论都是如此。我们将使用标量场论的语言，但论证可以推广到旋量和矢量场论。

首先考虑在具有任意数量场 $\phi^i$ 的场论中，对常数（与 $x$ 无关的）经典场计算有效势的问题。有效势的质量量纲为 4，因此我们预期 $V_{eff}(\phi_{cl})$ 会有直到 $\Lambda^4$ 的发散项。为理解这些发散，把 $V_{eff}(\phi_{cl})$ 作泰勒级数展开：
\begin{equation}
  V_{eff}(\phi_{cl}) = A_0 + A_2\phi_{cl}^i\phi_{cl}^i
    + \frac{1}{4!}A_4(\phi_{cl}^i\phi_{cl}^i)^2 + \cdots .
\end{equation}
（在没有对称性 $\phi^i\to-\phi^i$ 的理论中，可能还会有线性于和三次于 $\phi_{cl}$ 的项；为简单起见我们省略这些项。）系数 $A_0$、$A_2$、$A_4$ 的质量量纲分别为 4、2 和 0；因此我们预期它们分别含有 $\Lambda^4$、$\Lambda^2$ 和 $\log\Lambda$ 发散。幂次计数分析预言，泰勒级数展开中所有更高的项应当是有限的。

常数项 $A_0$ 与 $\phi_{cl}$ 无关；它没有物理意义。然而，$A_2$ 和 $A_4$ 中的发散出现在物理量中，因为这些系数进入逆传播子 \eqref{eq:11.90} 和不可约四点函数 \eqref{eq:11.96}，因此出现在 $S$ 矩阵元的计算中。有效作用量中还有一个系数按幂次计数具有非负质量量纲；这就是二次于 $\partial_\mu\phi_{cl}$ 的项的系数，它出现在对非常数背景场求值有效作用量时：
\begin{equation}
  \Delta\Gamma[\phi_{cl}] = \int d^4x\ B\ \frac{1}{2}(\partial_\mu\phi_{cl})^2 .
  \label{eq:11.98}
\end{equation}
有效作用量按 $\phi_{cl}$ 幂次泰勒展开中的所有其他系数按幂次计数都是有限的。

我们现在可以论证，原始拉格朗日量的抵消项足以消除计算 $\Gamma[\phi_{cl}]$ 时可能出现的发散。论证分两步进行。我们首先使用 BPHZ 定理论证，格林函数的发散可以通过调整一组抵消项来消除，这些抵消项对应于可以以非负质量量纲的系数添加到拉格朗日量中的可能算符。这些抵消项的系数与有效作用量的系数 $A_2$、$A_4$ 和 $B_2$ 一一对应。接下来，我们利用有效作用量显然对模型的原始对称群不变这一事实。即使模型的真空态具有自发对称性破缺，这也是成立的。有效作用量的这一对称性来自第~11.4 节的分析，因为我们在那里提出的计算有效作用量的方法显然对拉格朗日量的原始对称性不变。

把这两个结果结合起来，我们得出结论：即使原始对称性被自发破缺，也总是可以通过调整对理论原始对称性不变的那组抵消项来使有效作用量有限。利用第~11.5 节的结果（该节解释了如何从有效作用量的泛函导数构造理论的格林函数），这一可重整性结论扩展到理论的所有格林函数。

为使这一抽象论证更具体，我们将在简单例子中演示有效作用量的泛函导数如何产生一组其发散对应于对称抵消项的费曼图。那么，让我们再次回到 $O(N)$ 不变的线性西格马模型，并计算 $\Gamma[\phi_{cl}]$ 的二阶泛函导数。如果我们构造的整个形式体系自洽，我们就应该能够把结果识别为逆传播子的费曼图展开，其发散对应于 $O(N)$ 对称标量场论的抵消项。

首先，我们为这个模型显式写出表达式 \eqref{eq:11.63}：
\begin{equation}
  \Gamma[\phi_{cl}] = \int d^4x\ \Bigl( \tfrac{1}{2}(\partial_\mu\phi_{cl})^2
    + \tfrac{1}{2}m^2(\phi_{cl})^2 - \frac{\lambda}{4}\bigl[ (\phi_{cl})^2
      \bigr]^2 \Bigr) + \frac{i}{2}\log\det[-iD^{ij}] + \cdots ,
  \label{eq:11.99}
\end{equation}
其中
\begin{equation}
  -iD^{ij} = \partial^2\delta^{ij} + \bigl( \lambda(\phi_{cl}(x))^2 - \mu^2
    \bigr)\delta^{ij} + 2\lambda\phi^i_{cl}(x)\phi^j_{cl}(x) .
  \label{eq:11.100}
\end{equation}
对于常数 $\phi_{cl}$，$D^{ij}$ 是这样一个算符：它作用在标量场的给定分量上，等于质量平方由式~\eqref{eq:11.69} 给出的克莱因--戈登算符。这是线性西格马模型逆传播子的领头阶近似。

为找到逆传播子的高阶修正，我们必须计算 $\Gamma[\phi_{cl}]$ 中量子修正项的二阶泛函导数。从 \eqref{eq:11.99} 我们发现
\begin{equation}
  \begin{aligned}
    \frac{\delta^2\Gamma}{\delta\phi^i_{cl}(x)\delta\phi^j_{cl}(y)}
    &= iD^{ij}\delta(x-y) + \frac{\delta^2}{\delta\phi^i_{cl}(x)
      \delta\phi^j_{cl}(y)}\ \frac{i}{2}\log\det[-iD^{ij}] + \cdots .
  \end{aligned}
\end{equation}
第一项正是克莱因--戈登算符 $iD^{ij}\delta(x-y)$。为计算第二项，对矩阵行列式使用恒等式 \eqref{eq:9.78}：
\begin{equation}
  \frac{\partial}{\partial a}\log\det M(a) = \frac{\partial}{\partial a}
    \mathrm{tr}\log M(a) = \mathrm{tr}\Bigl( M^{-1}\frac{\partial M}
      {\partial a} \Bigr) .
  \label{eq:11.101}
\end{equation}
利用这个恒等式，我们发现
\begin{equation}
  \frac{\delta}{\delta\phi^k_{cl}(z)}\ \frac{i}{2}\log\det[-iD]
  = i\lambda\bigl( \phi^k_{cl}(z)\delta^{ij} + \phi^i_{cl}(z)\delta^{jk}
    + \phi^j_{cl}(z)\delta^{ik} \bigr)\ (D^{-1})^{ij}(z,z) .
  \label{eq:11.102}
\end{equation}
量 $(D^{-1})^{ij}(x,y)$ 是克莱因--戈登传播子。为再微分一次，我们可以使用恒等式式~(11.94)；这给出
\begin{equation}
  \begin{aligned}
    &\frac{\delta^2}{\delta\phi^k_{cl}(z)\delta\phi^\ell_{cl}(w)}\
    \frac{i}{2}\log\det[-iD]\\
    &= -\lambda\bigl( \delta^{k\ell}\delta^{ij} + \delta^{ik}\delta^{j\ell}
      + \delta^{i\ell}\delta^{jk} \bigr)\ (D^{-1})^{ij}(z,z)\delta(z-w)\\
    &\quad + 2i\lambda^2\bigl( \phi^k_{cl}(z)\delta^{im} + \phi^i_{cl}(z)
      \delta^{mk} + \phi^m_{cl}(z)\delta^{ik} \bigr)\ (D^{-1})^{im}(z,w)\
      \bigl( \phi^\ell_{cl}(w)\delta^{nj} + \phi^n_{cl}(w)\delta^{n\ell}
        + \phi^{n}_{cl}(w)\delta^{\ell n} \bigr)\ (D^{-1})^{nj}(w,z) .
  \end{aligned}
  \label{eq:11.103}
\end{equation}
这预期是逆传播子在单圈阶的形式修正，而且我们确实可以在 \eqref{eq:11.103} 中识别出单圈图的值。注意，在这一推导中，对 $D^{-1}$ 的每一次泛函导数都给图增加一个传播子，从而降低发散度，这与我们在第~10.1 节的一般论证一致。

这个例子说明，$\Gamma[\phi_{cl}]$ 的各阶泛函导数是通过费曼图展开来计算的，其中的传播子和顶点依赖于经典场。当经典场是常数时，传播子约化为普通的克莱因--戈登传播子，因此 BPHZ 定理适用。所有紫外发散都可以通过使用最一般的质量、顶点和场强重整化集合，从对 $\Gamma[\phi_{cl}]$ 求导得到的所有振幅中消除。与此同时，微扰论显然对原始拉格朗日量的对称性不变，因此出现的唯一发散——从而所需的唯一抵消项——就是尊重这一对称性的那些。于是，一般而言，在对称群下不变的标量场可重整化理论的所有振幅，都可以只使用对该对称性不变的抵消项集合而变得有限。这为第~11.2 节开头提出的问题给出了一个完整且相当令人满意的回答。

在空间变化背景场中计算有效作用量，尚未达到 BPHZ 定理证明所涉及的严格程度。然而，可以预期在这种情况下，对称理论的标准抵消项集合也应当充分。我们可以直觉地论证这一点，利用费曼图的紫外发散在时空中是局域的这一事实。因此，为了理解在平滑变化的背景 $\phi_{cl}(x)$ 中计算的发散，我们可以把时空分成小盒子，在每个盒子中 $\phi_{cl}(x)$ 近似为常数，并按导数 $\partial_\mu\phi_{cl}(x)$ 展开。在这个按 $\phi_{cl}(x)$ 幂次的展开中，泰勒级数系数是常数背景下的 $\Gamma$ 泛函导数，而我们已知它们是可以重整化的。这一结论扩展到所有阶，对称理论的抵消项对于具有自发对称性破缺的理论之重整化是充分的。

\subsection{再访戈德斯通定理}

作为有效作用量形式体系的一个最终应用，让我们回到戈德斯通定理在存在量子修正时是否仍然成立的问题。回忆我们在第~11.1 节末尾在经典层面证明了这一定理：我们在 \eqref{eq:11.13} 中证明，如果拉格朗日量有一个被自发破缺的连续对称性，那么经典势 $V(\phi)$ 的二阶导数矩阵有一个相应的零特征值。根据式~\eqref{eq:11.11}，这意味着经典理论包含一个无质量标量粒子，它与被自发破缺的对称性相联系。

利用有效作用量形式体系，这一论证几乎可以逐字地在完整量子场论中重复。有效势 $V_{eff}(\phi_{cl})$ 概括了理论的完整解，包括所有阶的量子修正。同时，它满足经典势的一般性质：它对理论的对称性不变，其极小值给出 $\phi$ 的真空期望值。这意味着我们在 \eqref{eq:11.13} 中给出的论证对 $V_{eff}$ 和对 $V$ 以完全相同的方式成立：如果原始拉格朗日量的一个连续对称性被 $\langle\phi\rangle$ 自发破缺，那么 $V_{eff}(\phi_{cl})$ 的二阶导数矩阵在对称方向上有一个零特征值。

我们现在论证，正如在经典层面一样，这样一个零特征值的出现意味着无质量标量粒子的存在。在我们关于有效作用量一般性质的讨论中，我们证明了它的二阶泛函导数是逆传播子，并且通过式~\eqref{eq:11.97}，这一导数给出量子理论中的质量谱。让我们为包含多个标量场的理论重写式~\eqref{eq:11.97}：
\begin{equation}
  \int d^4x\ e^{ip\cdot(x-y)}\ \frac{\delta^2\Gamma[\phi_{cl}]}
    {\delta\phi^i_{cl}(x)\delta\phi^j_{cl}(y)} = 0 .
  \label{eq:11.105}
\end{equation}
质量为 $m$ 的粒子对应于这个矩阵方程在 $p^2 = m^2$ 处的零特征值。现在令 $p = 0$。这意味着我们对常数场微分 $\Gamma[\phi_{cl}]$。因此，我们可以用 $\Gamma[\phi_{cl}]$ 在常数经典场下的值来替换它，而后者正是有效势。我们发现，当二阶导数矩阵
\begin{equation}
  \frac{\partial^2 V_{eff}}{\partial\phi^i_{cl}\partial\phi^j_{cl}} ,
\end{equation}
有一个零特征值时，量子场论包含一个零质量标量粒子。这就完成了戈德斯通定理的证明。

这一关于戈德斯通定理的论证说明了有效作用量形式体系的威力。该形式体系给出了一个对任何阶量子修正都有效的自发对称性破缺的几何图像。作为额外收获，它由以简单方式重整化的对象构建而成。这一形式体系在理解自发对称性破缺的应用时将被证明是有用的，这些应用在本书余下部分中以几种不同情境出现。

\problems

\textbf{11.1} \emph{自旋波理论。}
\begin{enumerate}[label=(\alph*)]
  \item 证明下面这个奇妙的公式：设 $\phi(x)$ 是一个自由标量场，其传播子为 $\langle T\phi(x)\phi(0)\rangle = D(x)$。则
    \begin{equation}
      \langle T\,e^{i\phi(x)}e^{-i\phi(0)} \rangle
      = e^{\left[ D(x) - D(0) \right]} .
    \end{equation}
    （因子 $D(0)$ 给出对整体归一化的一个形式上发散的调整。）
  \item 我们可以在欧几里得场论中使用这个公式来讨论 $T < T_c$ 时具有自发对称性破缺理论中的关联函数。让我们只考虑破缺的 $O(2)$ 或 $U(1)$ 对称性的最简单情形。我们可以把局域自旋密度写成一个复变量
    \begin{equation}
      s(x) = s_1(x) + is_2(x) .
    \end{equation}
    全局对称性是变换
    \begin{equation}
      s(x) \to e^{-i\alpha}s(x) .
    \end{equation}
    如果我们假设物理把 $s(x)$ 的模冻结，我们可以参数化
    \begin{equation}
      s(x) = Ae^{i\phi(x)}
    \end{equation}
    并为场 $\phi(x)$ 写一个有效拉格朗日量。理论的对称性变为平移对称性
    \begin{equation}
      \phi(x) \to \phi(x) - \alpha .
    \end{equation}
    证明（对于 $d > 0$）与这一对称性一致的最一般的可重整化拉格朗日量是自由场论
    \begin{equation}
      \mathcal{L} = \tfrac{1}{2}\rho(\nabla\phi)^2 .
    \end{equation}
    在统计力学中，常数 $\rho$ 称为\emph{自旋波模量}。对 $\rho$ 的一个合理假设是：它在 $T < T_c$ 时有限，并且当 $T$ 从下方趋于 $T_c$ 时趋于 0。
  \item 计算关联函数 $\langle s(x)s^{*}(0)\rangle$。调整 $A$ 以给出一个物理上合理的归一化（假设系统在单个原子间距的标度处有一个物理截断），并展示这一关联函数对 $x$ 的依赖对于 $d = 1, 2, 3, 4$ 的显式形式。解释你的结果的意义。
\end{enumerate}

\textbf{11.2} \emph{一个零阶自然关系。} 本题研究一个与费米子耦合的 $N = 2$ 线性西格马模型：
\begin{equation}
  \mathcal{L} = \tfrac{1}{2}(\partial_\mu\phi^i)^2 + \tfrac{1}{2}\mu^2(\phi^i)^2
    - \frac{\lambda}{4}\bigl[ (\phi^i)^2 \bigr]^2
    + \bar{\psi}(i\slashed{\partial})\psi
    - g\bar{\psi}(\phi^1 + i\gamma^5\phi^2)\psi ,
  \label{eq:11.p1}
\end{equation}
其中 $\phi^i$ 是一个二分量场，$i = 1, 2$。
\begin{enumerate}[label=(\alph*)]
  \item 证明该理论具有以下全局对称性：
    \begin{equation}
      \phi^1 \to \cos\alpha\,\phi^1 - \sin\alpha\,\phi^2,\qquad
      \phi^2 \to \sin\alpha\,\phi^1 + \cos\alpha\,\phi^2,\qquad
      \psi \to e^{i\gamma^5\alpha/2}\psi .
      \label{eq:11.p2}
    \end{equation}
    并证明经典运动方程的最小能量解会自发地破缺这一对称性。
  \item 用 $v$ 表示场 $\phi^1$ 的真空期望值，并作变量替换
    \begin{equation}
      \phi^1(x) = (v + \sigma(x), \pi(x)) .
      \label{eq:11.p3}
    \end{equation}
    用这些新变量写出拉格朗日量，并证明费米子获得一个由下式给出的质量
    \begin{equation}
      m_f = g\cdot v .
      \label{eq:11.p4}
    \end{equation}
  \item 计算对 $m_f$ 的单圈辐射修正，选择重整化条件使 $v$ 和 $g$（定义为零动量转移处的 $\bar{\psi}\psi\pi$ 顶点）不接受辐射修正。证明关系 \eqref{eq:11.p4} 接受非零修正，但这些修正是有限的。这与我们在第~11.6 节的一般讨论一致。
\end{enumerate}

\textbf{11.3} \emph{Gross--Neveu 模型。} Gross--Neveu 模型是二维时空中的一个费米子模型，具有离散手征对称性：
\begin{equation}
  \mathcal{L} = \bar{\psi}_i i\slashed{\partial}\psi^i
    + \tfrac{1}{2}g^2(\bar{\psi}_i\psi^i)^2 ,
\end{equation}
其中 $i = 1, \ldots, N$。二维费米子的动能项由满足二维狄拉克代数的矩阵 $\gamma^\mu$ 构建。这些矩阵可以是 $2\times2$ 的：
\begin{equation}
  \gamma^0 = \sigma^2,\qquad \gamma^1 = i\sigma^1 ,
\end{equation}
其中 $\sigma^i$ 是泡利西格马矩阵。定义
\begin{equation}
  \gamma^5 = \gamma^0\gamma^1 = \sigma^3 ;
\end{equation}
这个矩阵与 $\gamma^\mu$ 反对易。
\begin{enumerate}[label=(\alph*)]
  \item 证明该理论对
    \begin{equation}
      \psi^i \to \gamma^5\psi^i ,
    \end{equation}
    不变，并且这一对称性禁止费米子质量的出现。
  \item 证明该理论在 2 维中是（在量纲分析的层面上）可重整化的。
  \item 证明该理论的泛函积分可以表示成以下形式：
    \begin{equation}
      Z = \int \mathcal{D}\sigma\ \exp\Bigl[ -\frac{i}{2}\int d^2x\ \sigma^2
        \Bigr]\ \int \mathcal{D}\bar{\psi}\,\mathcal{D}\psi\ \exp\Bigl[
        i\int d^2x\ \bar{\psi}(i\slashed{\partial} + g\sigma)\psi \Bigr] ,
    \end{equation}
    其中 $\sigma(x)$（不要与泡利矩阵混淆）是一个没有动能项的新标量场。
  \item 通过对费米子场 $\psi$ 积分，计算对 $\sigma$ 有效势的领头修正。你将遇到一个狄拉克算符的行列式；为求值这一行列式，先把算符对角化：首先转到傅里叶分量，然后对角化与每个傅里叶模式相联系的 $2\times2$ 泡利矩阵。（或者，你也可以直接取这个 $2\times2$ 矩阵的行列式。）这个 1 圈贡献需要一个正比于 $\sigma^2$ 的重整化（即对 $g^2$ 的重整化）。用最小减除进行重整化。
  \item 忽略两圈及更高阶的贡献，把这个势极小化。证明 $\sigma$ 场获得一个真空期望值，它破缺了第 (a) 部分的对称性。说服你自己这一结果不依赖于所选择的特定重整化条件。
  \item 注意第 (e) 部分导出的有效势按如下形式依赖于 $g$ 和 $N$
    \begin{equation}
      V_{eff}(\sigma_{cl}) = N\cdot f(g^2N) .
    \end{equation}
    （整体的 $N$ 因子在一个含 $N$ 个场的理论中是可以预期的。）构造有效势的几个高阶贡献，并证明它们含有额外的 $N^{-1}$ 因子，如果我们取极限 $N\to\infty$（保持 $(g^2N)$ 固定），这些因子会抑制它们。在这个极限下，第 (e) 部分的结果是无歧义的。
\end{enumerate}
